% Part: lambda-calculus
% Chapter: introduction
% Section: syntax

\documentclass[../../../include/open-logic-section]{subfiles}

\begin{document}

\olfileid{lam}{int}{syn}
\olsection{نحو حساب لامبدا}

نبدأ بمتتالية من المتغيرات~$x$ و$y$ و$z$،~\dots وببعض رموز الثوابت
$a$ و$b$ و$c$،~\dots. وتعرّف مجموعة الحدود استقرائيًا كما يأتي:
\begin{enumerate}
\item كل متغير حد.
\item كل ثابت حد.
\item إذا كان~$M$ و$N$ حدين، فإن~$(MN)$ حد أيضًا.
\item إذا كان~$M$ حدًا وكان~$x$ متغيرًا، فإن~$(\lambd[x][M])$ حد.
\end{enumerate}
تسمى الحدود من الصورة~$(MN)$ \emph{تطبيقات}، وتسمى الحدود من الصورة
$(\lambd[x][M])$ \emph{تجريدات}.

يسمى النسق الخالي تمامًا من الثوابت حساب لامبدا \emph{الخالص}. وسنعمل
أساسًا في حساب~$\lambd$ الخالص، ولذلك ترمز جميع الحروف الصغيرة إلى
متغيرات. ونستعمل الحروف الكبيرة ($M$ و$N$، إلخ) لترمز إلى حدود حساب
$\lambd$.

سنتبع بعض الاصطلاحات الترميزية:

\begin{conv}
\begin{enumerate}
\item حين تحذف الأقواس، يجري التطبيق من اليسار إلى اليمين. فمثلًا، إذا
  كانت~$M$ و$N$ و$P$ و$Q$ حدودًا، فإن~$MNPQ$ اختصار لـ$(((MN)P)Q)$.
\item وكذلك، حين تحذف الأقواس، يعطى تجريد لامبدا أوسع مجال ممكن. فمثلًا،
  تقرأ~$\lambd[x][MNP]$ على أنها~$(\lambd[x][((MN)P)])$.
\item يمكن استعمال لامبدا لتجريد متغيرات متعددة. فمثلًا، تختصر
  $\lambd[xyz][M]$ التعبير
  $\lambd[x][\lambd[y][\lambd[z][M]]]$.
\end{enumerate}
\end{conv}

فمثلًا، يختصر
\[
\lambd[xy][xxyx \lambd[z][xz]]
\]
التعبير
\[
\lambd[x][\lambd[y][((((xx)y)x)(\lambd[z][(xz)]))]].
\]
ينبغي أن تحفظ هذه الاصطلاحات. ستربكك في البداية، لكنك ستعتادها، وبعد
مدة ستجدها أقل إرباكًا من الاضطرار إلى معالجة غابة من الأقواس.

يسمى حدان لا يختلفان إلا في أسماء المتغيرات المربوطة متكافئين
$\alpha$؛ ومن ذلك~$\lambd[x][x]$ و$\lambd[y][y]$. ومن الملائم أن نعدهما
الحد ``نفسه''؛ أي إننا حين نقول إن~$M$ و$N$ متماثلان، نعني أيضًا
``حتى إعادة تسمية المتغيرات المربوطة''. وتسمى المتغيرات الواقعة في مجال
$\lambd$ ``مربوطة''، وتسمى سائر المتغيرات ``حرة''. ولا توجد متغيرات حرة
في المثال السابق، لكن في
\[
(\lambd[z][yz])x
\]
يكون~$y$ و$x$ حرين، ويكون~$z$ مربوطًا.

\end{document}
